\input{../tex-inputs/latex-korrekturansicht-vorspann}

\section[Arthur Schnitzler an Rainer Maria Rilke, 20. 1. 1900]{L04342 Arthur Schnitzler an Rainer Maria Rilke, 20. 1. 1900}
\nopagebreak\mylabel{L04342v}
\rehead{ }\normalsize\beginnumbering\briefempfaengerindex{Rilke, Rainer Maria@\textsc{Rilke, Rainer Maria}!zzzSchnitzler, Arthur@\emph{von Arthur Schnitzler}!1900-01-201@{20. 1. 1900}|(be}
\toendnotes[C]{\smallbreak\pagebreak[2]}
\correspDesc{Versand  durch Arthur Schnitzler am 20. 1. 1900 in Wien
\newline{}Erhalt  durch Rainer Maria Rilke im Zeitraum [21. 1. 1900 – 25. 1. 1900?] in Berlin}\toendnotes[C]{\smallbreak}
\Standort{DLA, A:Schnitzler, HS.2022.81.}
\physDesc{Briefkarte, 67 Zeichen
\newline{}Handschrift: schwarze Tinte, deutsche Kurrent}\toendnotes[C]{\smallbreak}
\pstart
           \noindent{}{\pb}Herzlichen Dank für das ſchöne \label{K_L04342-1v}\edtext{\textcolor{green}{Buch}\pwindex{Rilke, Rainer Maria 4.\,12.\,1875 Prag – 29.\,12.\,1926 Montreux@\textsc{Rilke, Rainer Maria} (4.\,12.\,1875 Prag – 29.\,12.\,1926 Montreux), \emph{Schriftsteller}!Mir zur Feier@\strich\emph{Mir zur Feier}|pwuv}{}\ledrightnote{{$\rightarrow$}\emph{\textcolor{green}{Mir zur Feier}}}}{\lemma{\textnormal{\emph{Buch}}}\Cendnote{\textnormal{Vermutlich:
                        \textcolor{blue}{Rainer Maria Rilke}\pwindex{Rilke, Rainer Maria 4.\,12.\,1875 Prag – 29.\,12.\,1926 Montreux@\textsc{Rilke, Rainer Maria} (4.\,12.\,1875 Prag – 29.\,12.\,1926 Montreux), \emph{Schriftsteller}|pwk}: \emph{\textcolor{green}{Mir zur Feier. Gedichte}\pwindex{Rilke, Rainer Maria 4.\,12.\,1875 Prag – 29.\,12.\,1926 Montreux@\textsc{Rilke, Rainer Maria} (4.\,12.\,1875 Prag – 29.\,12.\,1926 Montreux), \emph{Schriftsteller}!Mir zur Feier@\strich\emph{Mir zur Feier}|pwk}}.
                     Berlin: \emph{\textcolor{brown}{Georg Heinrich
                        Meyer}\orgindex{Georg Heinrich Meyer@Georg Heinrich Meyer|pwk}}{ }1899. Das
                  Buch erschien Ende 1899 und wurde Anfang des Jahres verteilt, vgl.
                  die Widmung eines Exemplars an \textcolor{blue}{Raphael
                     Löwenfeld}\pwindex{Löwenfeld, Raphael 11.\,2.\,1854 Poznan – 28.\,12.\,1910 Berlin@\textsc{Löwenfeld, Raphael} (11.\,2.\,1854 Poznan – 28.\,12.\,1910 Berlin), \emph{Theaterleiter}|pwk}: »Herrn Direktor Dr. \textsc{\textcolor{blue}{Raphael Löwenfeld}\pwindex{Löwenfeld, Raphael 11.\,2.\,1854 Poznan – 28.\,12.\,1910 Berlin@\textsc{Löwenfeld, Raphael} (11.\,2.\,1854 Poznan – 28.\,12.\,1910 Berlin), \emph{Theaterleiter}|pw}} als
                        verſpäteten Dank für das liebe ruſſiſche Geleite und alle damit
                        zuſammenhängende Güte!{ / }Rainer Maria Rilke.{ / }\textcolor{pink}{Schmargendorf bei Berlin}\oindex{Schmargendorf@\textbf{Schmargendorf}, \emph{Ehemaliger Ort}|pw},
                           anfangs 1900.« (\url{https://doi.org/10.7891/e-manuscripta-53688}) }}}\label{K_L04342-1}.\pend
           \pstart Ihr \spacefill\mbox{Arthur Schnitzler}\pend{}
\pstart
           20. 1. 1900.\pend
           \selectlanguage{ngerman}\endnumbering\briefempfaengerindex{Rilke, Rainer Maria@\textsc{Rilke, Rainer Maria}!zzzSchnitzler, Arthur@\emph{von Arthur Schnitzler}!1900-01-201@{20. 1. 1900}|)be}\mylabel{L04342h}
\begin{anhang}
\end{anhang}\newcommand{\dateiname}{L04342}\newcommand{\titel}{Arthur Schnitzler an Rainer Maria Rilke, 20. 1. 1900}\newcommand{\editorInnen}{Herausgegeben von Jahnke, SelmaMüller, Martin Anton}\input{../tex-inputs/latex-korrekturansicht-abspann}
